Particle Swarm Optimization (PSO) remains one of the most widely adopted metaheuristic algorithms for continuous optimization, yet its tendency to converge prematurely on complex landscapes motivates the exploration of quantum-inspired variants. This paper presents a comprehensive empirical comparison between Classical PSO and Quantum-behaved PSO (QPSO) using a modular, open-source Python framework that supports three distinct quantum randomness generation modes: a lightweight mathematical simulation, full Qiskit-based quantum circuit execution, and a hybrid approach that balances computational cost with randomness quality. We evaluate both algorithms on six standard benchmark functions---Sphere, Rastrigin, Rosenbrock, Ackley, Griewank, and Schwefel---spanning unimodal, multimodal, and valley-shaped landscapes across dimensionalities ranging from 2 to 20. Each experimental configuration is assessed over 50 independent runs with paired random seeds, and results are validated using the Wilcoxon signed-rank test at a significance level of $\alpha = 0.05$. Rather than a uniform advantage for either method, our findings reveal a \emph{dimensional crossover}: Classical PSO is significantly more accurate on smooth and low-dimensional landscapes, where its velocity-based exploitation reaches near machine precision, whereas QPSO overtakes it at high dimensionality ($n = 20$) on the Ackley, Sphere, and deceptive Schwefel functions, where classical PSO converges prematurely. Across the 24 benchmark--dimensionality combinations, Classical PSO wins 9, QPSO wins 6, and 9 show no significant difference. On the archetypal multimodal Rastrigin function the mean-fitness difference is not significant, yet a binary local-minimum escape-rate analysis shows QPSO reaching the global basin significantly more often at $n = 5$ (41/50 vs.\ 27/50, Fisher's exact $p = 0.005$), a direct demonstration of the quantum-tunnelling behaviour that motivates the method. The complete experimental framework---which additionally supports full and hybrid Qiskit-based randomness modes for future study---is released as an open-source package with reproducible pipelines for convergence analysis, statistical testing, and visualization.
